\section{第17组：原子核与核物理}
\sourceproblem{17-T1}{106}{110}
\problemtext 钴60源每秒测得3700个伽马光子，测量效率10\%，半衰期5.27年，求活度与质量。
\workings 令每次衰变计入测量的平均光子数为 $Y$，光子探测效率 $\eta=0.10$，则 $A=3700/(\eta Y)=3.70\times10^4/Y$ Bq。$N=AT_{1/2}/\ln2$，$m=NM/N_A$，取 $M=0.060$ kg/mol，得到 $m\approx8.85\times10^{-13}/Y$ kg。教材若取一次衰变计一个光子或效率已按每次衰变定义，则 $Y=1$；若两个级联光子均计数且各自效率10\%，近似 $Y=2$，活度和质量均减半。
\crosscheck 活度是核衰变次数率，不是光子计数率；原题未说明能窗和光子产额，保留条件解。
\sourceproblem{17-T2}{106}{110}
\problemtext 10 kg的天然铀氧化物 $U_3O_8$ 中估计镭226、镤231含量，给定铀同位素丰度和两子核半衰期。
\workings 取化学式纯物质模型，铀原子总数 $N_U=3(10\,\mathrm{kg})N_A/M_{U_3O_8}$，用原子丰度 $f_{238}=0.9927,f_{235}=0.0072$。若矿物封闭且长衰变链达到长期平衡，则子核与相应母核活度相等，
\[
N_{\rm Ra}=f_{238}N_U T_{\rm Ra}/T_{238},\quad
N_{\rm Pa}=f_{235}N_U T_{\rm Pa}/T_{235}.
\]
质量分别为 $N_{\rm Ra}M_{\rm Ra}/N_A,N_{\rm Pa}M_{\rm Pa}/N_A$。题给子核半衰期1600年和 $3.27\times10^4$ 年，母核半衰期未在本题列出，故保留参数式。
\crosscheck 实际矿石的纯度、年龄、化学迁移及衰变链平衡状态影响含量；未达平衡时仅凭铀丰度不能确定子核存量。
\sourceproblem{17-T3}{107}{111}
\problemtext 古木片燃烧得0.3 mol二氧化碳，碳14活度读数每分钟9次，估算年龄。
\workings 还原碳质量约3.6 g。设活体参考比活度 $a_0$（每克每分钟衰变次数），效率 $\eta$，半衰期 $T_{1/2}$，则 $9=\eta(3.6)a_0\,2^{-t/T_{1/2}}$，故
\[
t=\frac{T_{1/2}}{\ln2}\ln\frac{3.6\eta a_0}{9}.
\]
若读数已校正为真实活度、取 $\eta=1$，仍需参考比活度与半衰期才能代入。
\crosscheck 原题没有给初始参考活度和计数效率；实际年代还需本底、同位素分馏及校准曲线等修正，不能凭给定读数唯一确定年龄。
\sourceproblem{17-T4}{107}{111}
\problemtext 太阳含氢 $1.5\times10^{30}$ kg，光度 $3.9\times10^{26}$ W，以恒功率估算氢燃尽时间。
\workings 设可参与聚变的氢比例为f，转为所计辐射能的质量比例为 $\epsilon$，则 $t=f\epsilon Mc^2/P$。取教材理想全部氢燃烧 $f=1$、氢聚变为氦的近似质量亏损 $\epsilon=0.007$，得到 $t\approx2.42\times10^{18}$ s，即约 $7.67\times10^{10}$ 年。
\crosscheck 不可把全部氢静能直接当作辐射能；真实太阳只有部分氢能参与主序阶段核心聚变，光度也会变化，故本数值不是太阳真实剩余寿命。
